% 原书分章PDF第25—28页：第26.5—26.6节。
\section{结语}

如果目标参数具有高阶路径可微性，那么，使用适当高阶梯度实现高阶TMLE，似乎是值得采用的做法。与一阶方法相比，高阶TMLE能够在较弱条件下成为渐近有效的代入估计量。本章给出了构造二阶TMLE的模板；任意阶的高阶TMLE也可以用类似方式实现。其实现与普通TMLE相同，只是要使用维数更高的最不利参数子模型，生成由各阶梯度确定的一组得分。

对于不具高阶路径可微性的目标参数，本章使用近似高阶梯度。这一近似策略，仅影响通过高阶梯度实现的高阶偏差减少。由于这样的高阶TMLE经过专门构造，能够消除TMLE渐近展开中的高阶项，因此，只要高阶梯度得到了合理近似，就可以期待实践中相对于标准一阶TMLE有实质性的偏差减少。遗憾的是，这些近似依赖可能为高维的平滑，于是人们可能怀疑，我们是否只是把一个难题换成了另一个难题。所幸，基于近似的高阶TMLE仍处于一个安全范围内，一阶TMLE的理想性质会得到保留。在最不利子模型中加入额外参数，总会在渐近意义上改善估计量，而且我们有理由预期，它不会实质性损害估计量的有限样本表现。

实践中，无论基于真实还是近似的高阶部分典范梯度，高阶TMLE都应以通过高阶最不利子模型尽可能多地减少偏差为目标。本章为实现这一点，提供了理论可靠且切实可行的构件。由于任何高阶TMLE都是代入估计量，拟合高阶最不利子模型所引入的有限样本变异，会受到模型及目标参数映射的全局界的控制。对于高阶程序，这一点可能更加关键，因为需要修正的项更多。更为重要的是，看似艰巨的调节参数及算法其他潜在输入的选择任务，可以借助C-TMLE的实现自然地解决。C-TMLE提供了具体工具，能够根据数据自适应地精细调整高阶最不利子模型的选择，以最大限度地发挥它减少偏差的作用。

正如本章所强调的，示例中的二阶TMLE具体展示了相对于一阶TMLE的改进：它能在明显更大的统计模型中渐近线性且有效。这与 Robins 等（2009）在一步估计量背景下的奠基性工作所取得的进展直接对应。像本章这样，将高阶部分典范梯度纳入TMLE框架的另一项优势是：这些梯度会对高阶TMLE的影响函数作出贡献，而这些贡献可以直接用于置信区间的构造，从而可能改善有限样本表现。

最后，由于定理中施加的Donsker类条件，在某些情形下可能较为严格，因此，发展高阶交叉验证TMLE显然很有意义。后者将在迭代更新程序中使用交叉验证版本的经验风险；即使这些Donsker类条件不成立，也可以证明所得估计量渐近线性且有效。这将直接推广 Zheng 和 van der Laan（2010）在一阶情形下的工作，尤其因为，正如我们已经展示的，高阶TMLE可以表述为最不利子模型经过扩充的一阶TMLE。

\section{注记与延伸阅读}

这一推广的基本构件早在数十年前就已建立，尤其是在 J. Pfanzagl 的工作（Pfanzagl，1985）中；其中引入了高阶梯度的概念，并对任意模型空间上有限维参数的高阶展开作出了形式化处理。本章旨在利用高阶展开，使得在丰富的无限维模型中构造统计参数的正则渐近线性估计量成为可能，并为这种构造提供指导。这也正是 J. Robins、L. Li、E. Tchetgen 和 A. van der Vaart 的奠基性贡献所遵循的视角（例如 Robins 等，2008a，2009；Li 等，2011；van der Vaart，2014）；这些作者最早为精确处理这一问题提供了严格框架。最早的论述来自 Robins 等（2008a），主要关注在无法进行正则估计的情形下，利用高阶梯度推导最优估计量。后续工作则致力于发展一步估计量的高阶对应方法；一步估计量的早期介绍例如见 Levit（1975）、Ibragimov 和 Khasminskii（1981）、Pfanzagl（1982）及 Bickel（1982）。其一般方法是，识别朴素代入估计量渐近展开中的主导项，并据此对这一朴素估计量作加性修正。Robins 等（2009）仔细确立了将这一方法扩展至二阶所需的基础，并在两个问题中展示其用法：一是估计密度函数平方的积分，二是估计缺失数据模型中的平均响应。Li 等（2011）十分详细地研究了利用这一方法推断处理效应的问题，并在某些无法进行正则推断的情形中，确立了极小极大速率结果。

van der Vaart（2014）对一步估计量的一般高阶扩展作了出色综述。一步估计量简单而且易于描述。然而，在有限样本中，它的表现容易下降，因为它没有任何保障措施，来确保施加于朴素代入估计量的加性修正，不会使估计量接近、甚至超出参数空间的自然边界。一步方法在参数空间中作事后偏差修正；本章重点讨论的TMLE框架，则为构造底层数据生成分布的估计量，或计算感兴趣参数所需部分的估计量提供指导，使所得代入估计量具有正则且渐近线性的渐近行为。因此，修正步骤是在模型空间中完成的，而非参数空间中。早在这一方法正式发展之前数十年，Pfanzagl（1982）就已强调了构建这种方法的吸引力。由于所得估计量自动满足参数空间的界，它从不产生无意义的输出，例如超出单位区间的概率估计；而且在有限样本中，它可能优于那些渐近等价、却不具有代入形式的估计量。比较一阶TMLE与一步估计量时，就会遇到这一问题；但由于修正过程的复杂程度增加，它在高阶推断中可能更加重要。

正如 Robins 等（2009）所强调、并如本章所讨论的，对若干感兴趣的统计参数，高阶梯度并不存在。尽管如此，仍可以使用近似高阶梯度产生正则估计量。实践中，近似高阶梯度这一概念，必然涉及某些调节参数的选择。Robins 等（2009）讨论了这一要求，但既没有提供实际指导，这种指导似乎也并不容易建立。如后文所讨论，在这一背景下使用TMLE的一个优势在于：此类调节参数的选择，可以自然地嵌入协同TMLE的框架，下文将其简称为C-TMLE。因此，在基于高阶展开进行推断时，必须精细调参才能获得适当近似，这一要求所带来的实际复杂问题便能够方便地处理。此外，正如下文将讨论的，TMLE框架也提供了有用工具，力求保证所得高阶估计量的表现不会比通常的一阶TMLE更差。
